\documentclass{ctexbook}
\usepackage[
    backend=biber,
    style=chem-acs,
]{biblatex}
\addbibresource{ref.bib}
\usepackage{EC}
\begin{document}
\chapter{VIIA族元素}
\section{VIIA族元素综述}
VIIA族元素,即我们常说的\tbf{卤素},其研究历史由来已久,在化学中也有十分重要的地位. 1811年Schweigger研究氯的性质时提出了卤素这一概念,意味产生盐的元素.在汉语中,卤的本意是制盐时剩余的黑色母液.
\subsection{VIIA族元素的发现与分布}
\paragraph{氟}
在17世纪,人们就发现了具有腐蚀性的氢氟酸~\ce{HF}和具有荧光性质的萤石~\ce{CaF2}. 1812年, Ampere提出了氟元素的概念,其名字即来源于萤石.然而,受限于当时实验室的条件和安全措施,直到1886年氟单质才由Moissan制备得到. Moissan采用的办法是用~\ce{Pt/Ir}电极在~\ce{Pt}管内电解~\ce{KHF2}的无水~\ce{HF}溶液:
\begin{tightcenter}
    \ce{2KHF2 ->[\text{电解}] 2KF + H2 + F2}
\end{tightcenter}

\indent 卤素单质的活泼性质决定它们在自然界中几乎总是以游离态存在.地壳中~\ce{F}的丰度为第13位,与~\ce{Mn},~\ce{Ba}相近.主要的存在形式是:萤石~\ce{CaF2},冰晶石~\ce{Na3AlF6}和氟磷灰石~\ce{Ca5(PO4)3F}.其中只有萤石被广泛用于合成氟及相关的化合物;自然界中的冰晶石是一种稀有的矿物,仅在格陵兰岛存在有开采价值的矿床,工业中电解铝所需要的冰晶石是人工合成的;氟磷灰石则主要用于制备磷酸盐.\\
\indent 直至今日,大规模制备~\ce{F2}单质仍然采用Moissan的方法,即电解~\ce{KHF2}与~\ce{HF}的混合物,只是电解的温度维持在$\SI{100}{\celsius}$左右.
\paragraph{氯}
氯是卤素中最早被分离出的.1774年, Scheele用~\ce{MnO2}氧化~\ce{HCl}得到~\ce{Cl2}单质:
\begin{tightcenter}
    \ce{MnO2 + 4HCl ->[\Delta] MnCl2 + 2H2O + Cl2}
\end{tightcenter}
受燃素理论影响,他并没有意识到自己制备得到了一种新的单质.直到后来人们才认识到这种黄绿色的物质是~\ce{Cl2}单质.1811年, Davy根据其颜色将其命名为氯.\\
\indent 氯在地壳中的储存不多,主要以岩盐~\ce{NaCl}的形式存在.大部分氯元素以~\ce{Cl-}的形式在海水或内陆咸水湖中存在.\\
\indent 工业上大部分~\ce{Cl2}来源于氯碱工业,即电解~\ce{NaCl}溶液(通常直接由海水浓缩处理得到,因此大部分氯碱厂都设在海边):
\begin{tightcenter}
    \ce{2NaCl + 2H2O ->[\text{电解}] Cl2 + H2 + 2NaOH}
\end{tightcenter}
\paragraph{溴}
%\printbibliography[title={\Large 参考文献}]
\end{document}